\documentclass[10pt,a4paper]{amsart}
\usepackage[margin=19mm]{geometry}
\usepackage{amsmath,amssymb,mathtools}
\usepackage[scr=rsfs,cal=euler]{mathalfa}
\usepackage{xfrac}
\usepackage[unicode,hidelinks,bookmarks=false]{hyperref}
\newcommand{\ie}{i.e.\ }
\newcommand{\cf}{cf.\ }
\newcommand{\resp}{respectively\ }
\newcommand{\Subsection}{\subsection}
\providecommand{\nobreakdash}{\nobreak\hbox{-}}
\setlength{\emergencystretch}{2em}
\multlinegap0pt
\usepackage{mathtools}
\usepackage{amssymb}
\usepackage[shortlabels]{enumitem}
\newtheorem{lemma}{Lemma}
\newtheorem{proposition}{Proposition}
\newcommand{\FM}[1]{\Delta_{#1}}
\newcommand{\Mh}{\mathcal{M}_h}
\newcommand{\PLtrim}[2]{\mathcal{P}_{#1}^{-}\Lambda^{#2}}
\newcommand{\PLtrimz}[2]{\mathring{\mathcal P}_{#1}^{-}\Lambda^{#2}}
\newcommand{\Sh}{\mathcal{S}_h}
\newcommand{\cof}[1]{\operatorname{cof}_{#1}}
\newcommand{\trc}{\operatorname{tr}}
\newcommand{\vvvert}{\vert\kern-0.25ex\vert\kern-0.25ex\vert}
\title{Uniform discrete Poincar\'e inequalities for Hybrid High-Order differential forms on polyhedral meshes}
\author{Silvano Pitassi}
\date{Source: arXiv:2608.10713v3 (CC BY 4.0)}
\begin{document}
\maketitle

\noindent\emph{Source and licence.} Uniform discrete Poincar\'e inequalities for Hybrid High-Order differential forms on polyhedral meshes. Version arXiv:2608.10713v3 (CC BY 4.0), distributed by arXiv under the Creative Commons Attribution 4.0 International licence, http://creativecommons.org/licenses/by/4.0/, as recorded in the arXiv licence metadata for this exact version and shown on the abstract page https://arxiv.org/abs/2608.10713v3. Full licence notice: \texttt{licenses/arxiv-2608.10713-CC-BY-4.0.txt}.\par\smallskip

\noindent\textit{Editorial scope.} Counted unit: the article's proposition prop:functional.cochain.isomorphism (source0.tex lines 1945--1964). The article's complete original proof of it is delivered with it (source0.tex lines 1965--2203, one \texttt{proof} environment of 239 lines). No mathematics is written, repaired or re-derived: the statement and the proof are the article's own lines, in the article's own notation, and no retained line is substituted or re-flowed.  Retained uncounted as the source-local context the proof needs: lines 1443--1515; lines 1455--1466; lines 1469--1477; lines 1493--1502; lines 1504--1511; lines 1937--1943.  Nothing was omitted from any retained span, so the package carries no omission record.  No retained line contains a citation, so the wrapper carries no bibliography and no bibliography entry.  No retained line contains a figure environment, an image-inclusion command or drawing code, so no image is delivered and the package declares no asset dependency.  Editorial formatting only, in addition: an AMS \texttt{amsart} preamble that restates the article's own declarations and macros used by the retained lines, namely the environments \texttt{lemma}, \texttt{proposition} on separate counters, together with the standard packages those lines need.  The retained environments print in this document as Lemma 1, Proposition 1 in order of appearance, whereas the article prints the same items with the numbering of its own full document.  The complete native source of the article as distributed in the arXiv e-print for this exact version is bundled unchanged as \texttt{sources/207345/source0.tex}.
\begin{lemma}[Trace-compatible geometric decomposition]
\label{lem:geometric.extensions}
Let $j\in\{k+1,\ldots,n-1\}$, $q_j\coloneqq j-k-1$, and fix a polynomial degree $r\geq 1$.

For every pair of cells $s\subset f$ with $q_j\leq\dim s$, there exists a linear operator
\[
  \mathcal E_{s,f}^j\colon
  \PLtrimz{r}{q_j}(\Sh(s))
  \longrightarrow
  \PLtrim{r}{q_j}(\Sh(f))
\]
such that $\mathcal E_{f,f}^j=\mathrm{Id}$, and, for every $e\in\FM{\dim f-1}(\partial f)$,
\begin{equation}\label{eq:geometric.extension.trace}
  \trc_e
  \mathcal E_{s,f}^j\zeta_s
  =
  \begin{cases}
    \mathcal E_{s,e}^j\zeta_s,
    &s\subset e,
    \\[1mm]
    0,
    &s\not\subset e.
  \end{cases}
\end{equation}

Moreover, one has the direct geometric decomposition
\begin{equation}\label{eq:geometric.direct.sum}
\PLtrim{r}{q_j}(\Sh(f))
=\bigoplus_{\substack{
s\subset f\\
q_j\leq\dim s
}}
\mathcal E_{s,f}^j
\PLtrimz{r}{q_j}(\Sh(s)).
\end{equation}
Equivalently, every $\nu_f\in\PLtrim{r}{q_j}(\Sh(f))$ admits a unique representation
\[
  \nu_f
  =
  \sum_{\substack{
    s\subset f\\
    q_j\leq\dim s
  }}
  \mathcal E_{s,f}^j\zeta_s,
  \qquad
  \zeta_s\in
  \PLtrimz{r}{q_j}(\Sh(s)).
\]

Setting $d\coloneqq\dim f$, the following uniform stability estimate holds.
\begin{equation}\label{eq:geometric.norm.equivalence}
  \vvvert\nu_f\vvvert_f^2
  \simeq
  \sum_{\substack{
    s\subset f\\
    q_j\leq\dim s
  }}
  h_f^{d-\dim s}
  \vvvert\zeta_s\vvvert_s^2,
\end{equation}
In particular,
\begin{equation}\label{eq:geometric.extension.stability}
  \vvvert
  \mathcal E_{s,f}^j\zeta_s
  \vvvert_f
  \lesssim
  h_f^{(d-\dim s)/2}
  \vvvert\zeta_s\vvvert_s.
\end{equation}
%The hidden constants depend only on $n$, $r$, the regularity of
%the auxiliary simplicial submeshes, and the standing mesh-regularity
%parameters.
\end{lemma}

\begin{equation}
  \trc_e
  \mathcal E_{s,f}^j\zeta_s
  =
  \begin{cases}
    \mathcal E_{s,e}^j\zeta_s,
    &s\subset e,
    \\[1mm]
    0,
    &s\not\subset e.
  \end{cases}
\end{equation}

\begin{equation}
\PLtrim{r}{q_j}(\Sh(f))
=\bigoplus_{\substack{
s\subset f\\
q_j\leq\dim s
}}
\mathcal E_{s,f}^j
\PLtrimz{r}{q_j}(\Sh(s)).
\end{equation}

\begin{equation}
  \vvvert\nu_f\vvvert_f^2
  \simeq
  \sum_{\substack{
    s\subset f\\
    q_j\leq\dim s
  }}
  h_f^{d-\dim s}
  \vvvert\zeta_s\vvvert_s^2,
\end{equation}

\begin{equation}
  \vvvert
  \mathcal E_{s,f}^j\zeta_s
  \vvvert_f
  \lesssim
  h_f^{(d-\dim s)/2}
  \vvvert\zeta_s\vvvert_s.
\end{equation}

\begin{equation}\label{eq:geometric.coordinate.definition}
\left(\mathcal G_d^j\Psi^d\right)_{s,f}(q_s)
\coloneqq
\Psi_f^d\left(\mathcal E_{s,f}^jq_s\right)
\qquad
\forall q_s\in\PLtrimz{m}{q_j}(\Sh(s)).
\end{equation}

\begin{proposition}[Isomorphism with coface coefficient complexes]
\label{prop:functional.cochain.isomorphism}
Under the assumptions of Lemma~\ref{lem:geometric.extensions}, the map $\mathcal G_d^j$ is an isomorphism and
\begin{equation}\label{eq:cochain.commutativity}
\mathcal G_d^j\mathbb B_d^j
=
\left(
\bigoplus_{s\in\mathfrak S_h^j}\delta_{d,s}
\right)
\mathcal G_{d-1}^j.
\end{equation}
Moreover, $\mathcal G_d^j$ and $(\mathcal G_d^j)^{-1}$ are uniformly bounded so that
\begin{equation}\label{eq:coordinate.norm.equivalence}
\|\mathcal G_d^j\Psi^d\|_{
\bigoplus_{s\in\mathfrak S_h^j}\mathbb C_{d,s}^j
}
\simeq
\|\Psi^d\|_{\mathbb X_d^j}.
\end{equation}
\end{proposition}

\begin{proof}
By \eqref{eq:geometric.direct.sum}, every $\nu_f\in\PLtrim{m}{q_j}(\Sh(f))$ has a unique representation
\[
\nu_f
=
\sum_{\substack{
s\in\mathfrak S_h^j\\
s\subset f
}}
\mathcal E_{s,f}^jq_s,
\qquad
q_s\in\PLtrimz{m}{q_j}(\Sh(s)).
\]
Consequently, a functional on $\PLtrim{m}{q_j}(\Sh(f))$ is uniquely determined by its restrictions to the geometric components, and $\mathcal G_d^j$ in \eqref{eq:geometric.coordinate.definition} is an isomorphism.
Its inverse is characterised by
\[
\Psi_f^d
\left(
\sum_{\substack{
s\in\mathfrak S_h^j\\
s\subset f
}}
\mathcal E_{s,f}^jq_s
\right)
=
\sum_{\substack{
s\in\mathfrak S_h^j\\
s\subset f
}}
\left(
\mathcal G_d^j\Psi^d
\right)_{s,f}(q_s).
\]

To prove the commuting identity \eqref{eq:cochain.commutativity}, let $f\in\cof{d}(s)$ and $q_s\in\PLtrimz{m}{q_j}(\Sh(s))$.
Using \eqref{eq:geometric.extension.trace}, we obtain
\[
\begin{aligned}
\left(
\mathcal G_d^j\mathbb B_d^j\Psi^{d-1}
\right)_{s,f}(q_s)
&=
\sum_{e\in\FM{d-1}(\partial f)}
\epsilon_{f,e}
\Psi_e^{d-1}
\left(
\trc_e\mathcal E_{s,f}^jq_s
\right)
\\
&=
\sum_{\substack{
e\in\FM{d-1}(\partial f)\\
s\subset e
}}
\epsilon_{f,e}
\Psi_e^{d-1}
\left(
\mathcal E_{s,e}^jq_s
\right)
\\
&=
\left(
\delta_{d,s}\mathcal G_{d-1}^j\Psi^{d-1}
\right)_f(q_s).
\end{aligned}
\]
This proves \eqref{eq:cochain.commutativity}.

It remains to prove the norm equivalence in \eqref{eq:coordinate.norm.equivalence}.
Fix $f\in\FM{d}(\Mh)$ and, for every $s\in\mathfrak S_h^j$ such that $s\subset f$, set
\[
\lambda_{s,f}
\coloneqq
\left(
\mathcal G_d^j\Psi^d
\right)_{s,f}
\in
\left(
\PLtrimz{m}{q_j}(\Sh(s))
\right)'.
\]
We first prove the local equivalence
\begin{equation}\label{eq:local.coordinate.norm.equivalence}
\|\Psi_f^d\|_{(\PLtrim{m}{q_j}(\Sh(f)))'}^2
\simeq
\sum_{\substack{
s\in\mathfrak S_h^j\\
s\subset f
}}
h_f^{\dim s-d}
\|\lambda_{s,f}\|_
{\left(\PLtrimz{m}{q_j}(\Sh(s))\right)'}^2.
\end{equation}

Let $\nu_f\in\PLtrim{m}{q_j}(\Sh(f))$ and write its geometric decomposition as
\[
\nu_f
=
\sum_{\substack{
s\in\mathfrak S_h^j\\
s\subset f
}}
\mathcal E_{s,f}^jq_s,
\qquad
q_s\in\PLtrimz{m}{q_j}(\Sh(s)).
\]
By the definition of the coordinate functionals and a weighted Cauchy--Schwarz inequality,
\[
\begin{aligned}
|\Psi_f^d(\nu_f)|
&=
\left|
\sum_{\substack{
s\in\mathfrak S_h^j\\
s\subset f
}}
\lambda_{s,f}(q_s)
\right|
\\
&\le
\left(
\sum_{\substack{
s\in\mathfrak S_h^j\\
s\subset f
}}
h_f^{\dim s-d}
\|\lambda_{s,f}\|_
{\left(\PLtrimz{m}{q_j}(\Sh(s))\right)'}^2
\right)^{1/2}
\left(
\sum_{\substack{
s\in\mathfrak S_h^j\\
s\subset f
}}
h_f^{d-\dim s}
\vvvert q_s\vvvert_s^2
\right)^{1/2}.
\end{aligned}
\]
The stability of the geometric decomposition \eqref{eq:geometric.norm.equivalence} therefore gives
\[
\|\Psi_f^d\|_{(\PLtrim{m}{q_j}(\Sh(f)))'}^2
\lesssim
\sum_{\substack{
s\in\mathfrak S_h^j\\
s\subset f
}}
h_f^{\dim s-d}
\|\lambda_{s,f}\|_
{\left(\PLtrimz{m}{q_j}(\Sh(s))\right)'}^2.
\]

Conversely, for every admissible subcell $s\in\mathfrak S_h^j$ with $s\subset f$ and every $q_s\in\PLtrimz{m}{q_j}(\Sh(s))$,
\[
\begin{aligned}
|\lambda_{s,f}(q_s)|
&=
\left|
\Psi_f^d
\left(
\mathcal E_{s,f}^jq_s
\right)
\right|
\\
&\le
\|\Psi_f^d\|_{(\PLtrim{m}{q_j}(\Sh(f)))'}
\vvvert
\mathcal E_{s,f}^jq_s
\vvvert_f
\\
&\lesssim
h_f^{(d-\dim s)/2}
\|\Psi_f^d\|_{(\PLtrim{m}{q_j}(\Sh(f)))'}
\vvvert q_s\vvvert_s,
\end{aligned}
\]
where the last inequality follows from \eqref{eq:geometric.extension.stability}.
Hence
\[
h_f^{\dim s-d}
\|\lambda_{s,f}\|_
{\left(\PLtrimz{m}{q_j}(\Sh(s))\right)'}^2
\lesssim
\|\Psi_f^d\|_{(\PLtrim{m}{q_j}(\Sh(f)))'}^2.
\]
Summing over the subcells of $f$ and using their uniformly bounded number proves the reverse inequality in \eqref{eq:local.coordinate.norm.equivalence}.

We now sum the local equivalence over all $d$-cells.
By the definition of the product norms and by exchanging the order of summation,
\[
\begin{aligned}
\left\|
  \mathcal G_d^j\Psi^d
\right\|_{
  \bigoplus_{s\in\mathfrak S_h^j}
  \mathbb C_{d,s}^j
}^2
&=
\sum_{\substack{
  s\in\mathfrak S_h^j\\
  \dim s\leq d
}}
\sum_{f\in\cof{d}(s)}
h_f^{n-2d+\dim s}
\|\lambda_{s,f}\|_{
  \left(
    \PLtrimz{m}{q_j}(\Sh(s))
  \right)'
}^2
\\
&=
\sum_{f\in\FM{d}(\Mh)}
h_f^{n-d}
\sum_{\substack{
  s\in\mathfrak S_h^j\\
  s\subset f
}}
h_f^{\dim s-d}
\|\lambda_{s,f}\|_{
  \left(
    \PLtrimz{m}{q_j}(\Sh(s))
  \right)'
}^2
\\
&\simeq
\sum_{f\in\FM{d}(\Mh)}
h_f^{n-d}
\|\Psi_f^d\|_{
  \left(
    \PLtrim{m}{q_j}(\Sh(f))
  \right)'
}^2
\\
&=
\|\Psi^d\|_{\mathbb X_d^j}^2.
\end{aligned}
\]
This proves \eqref{eq:coordinate.norm.equivalence} and completes the proof.
\end{proof}

\end{document}
